\documentclass[11pt,a4paper]{article}
\usepackage[margin=1in]{geometry}
\usepackage{amsmath,amssymb,amsthm,mathtools}
\usepackage{hyperref}
\usepackage{enumitem}
\usepackage{microtype}
\setlength{\emergencystretch}{3em}

\newtheorem{theorem}{Theorem}
\newtheorem{lemma}[theorem]{Lemma}
\newtheorem{proposition}[theorem]{Proposition}
\newtheorem{corollary}[theorem]{Corollary}
\newtheorem{remark}[theorem]{Remark}
\theoremstyle{definition}
\newtheorem{definition}[theorem]{Definition}
\newtheorem{assumption}{Assumption}[section]

\DeclareMathOperator{\tr}{tr}
\DeclareMathOperator{\Cov}{Cov}
\DeclareMathOperator{\Var}{Var}
\DeclareMathOperator{\diag}{diag}
\newcommand{\R}{\mathbb{R}}
\newcommand{\E}{\mathbb{E}}
\newcommand{\N}{\mathcal{N}}
\newcommand{\cR}{\mathcal{R}}
\newcommand{\cRc}{[\cR(t)]^c}
\newcommand{\ohat}{\hat{\omega}_{\mathrm{TL}}^{(t)}}
\newcommand{\ostar}{\omega^{(t)*}}
\newcommand{\bhat}{\hat{\beta}^{(s)}}
\newcommand{\ohats}{\hat{\omega}_s^{(t)}}
\newcommand{\SigR}{\Sigma_{\cR(t)}}
\newcommand{\SigRc}{\Sigma_{\cRc}}
\newcommand{\SigRRc}{\Sigma_{\cR(t),\cRc}}
\newcommand{\SigcR}{\Sigma_{\cRc|\cR(t)}}

\title{Asymptotic Upper Bound for the Predictive Risk\\
of the Transfer Learning Estimator}
\author{}
\date{}

\begin{document}
\maketitle

%% ===================================================================
\section{Setup and Notation}\label{sec:setup}
%% ===================================================================

We consider the transfer learning model described in the problem.
For each site $m\in\{t,s\}$ the data satisfy
\[
  y^{(m)} = x^\top\beta^{(m)}+\varepsilon, \qquad
  x\sim\N(0,\Sigma),\quad
  \varepsilon\sim\N(0,\sigma^{2}),
\]
where $\beta^{(m)}\in\R^{p}$ is a random coefficient vector, and
$\varepsilon$ is independent of $x$.  We observe only
$x_{\cR(t)}\in\R^{p_t}$ at the target site, and construct the
transfer--learning estimator
\[
  \ohat
  =\arg\min_{\omega}\biggl[
  \frac{1}{n_t}\|Y^{(t)}-X_{\cR(t)}^{(t)}\omega\|_{2}^{2}
  +\lambda_\omega\|\omega-\hat\omega_{s}^{(t)}\|_{2}^{2}\biggr],
\]
with $\hat\omega_{s}^{(t)}=\Phi^{(t)}\bhat$ as defined in the
problem.  We write $n=n_t$, $p_t=|\cR(t)|$, and adopt the shorthand
\[
  X=X_{\cR(t)}^{(t)}\in\R^{n\times p_t},\quad
  Y=Y^{(t)}\in\R^{n},\quad
  \hat\Sigma=\tfrac{1}{n}X^{\top}X,\quad
  \lambda=\lambda_\omega,\quad
  \Phi=\Phi^{(t)}.
\]

%% ===================================================================
\section{Assumptions}\label{sec:assumptions}
%% ===================================================================

We collect all hypotheses used in the proof.

\begin{assumption}[Proportional high-dimensional limit]\label{ass:dim}
As $n\to\infty$, $p/n\to\gamma\in(0,\infty)$ and
$p_t/n\to\eta\in(0,\infty)$.
\end{assumption}

\begin{assumption}[Random regression coefficients]\label{ass:coef}
For each $m\in\{s,t\}$,
\[
  \E[\beta^{(m)}]=0,\qquad
  \Cov(\beta^{(m)})=\frac{\alpha_m^{2}}{p}\,I_p,
  \qquad \alpha_m\neq 0,
\]
and
\[
  \Cov(\beta^{(t)},\beta^{(s)})
  =\frac{\rho\,\alpha_t\alpha_s}{p}\,I_p,
  \qquad \rho\in(0,1].
\]
We assume $(\beta^{(t)},\beta^{(s)})$ is jointly Gaussian.
\end{assumption}

\begin{assumption}[Source estimation error]\label{ass:source}
$\lambda_{\max}\!\bigl(\E[(\bhat-\beta^{(s)})
(\bhat-\beta^{(s)})^{\top}]\bigr)\le C^{U}/p$
uniformly as $(n,p)\to\infty$.
\end{assumption}

\begin{assumption}[Independence of source and target data]\label{ass:indep}
The source estimator $\bhat$ is constructed from source-site data
that is independent of the target-site observations
$(X,\varepsilon^{(t)})$.  (Dependence through the jointly random
$\beta^{(s)},\beta^{(t)}$ is permitted.)
\end{assumption}

\begin{assumption}[Isotropic covariance]\label{ass:iso}
$\Sigma=I_p$.  (This simplification is discussed in
Remark~\ref{rem:general_Sigma}.)
\end{assumption}

\begin{assumption}[Well-separated features]\label{ass:features}
The index set $\cR(t)\subseteq[p]$ is non-random, with
$|\cR(t)|=p_t$.
\end{assumption}

\begin{assumption}[Bounded regularisation parameter]\label{ass:lambda}
$\lambda>0$ is a deterministic constant (independent of the data)
that stays bounded away from $0$ and $\infty$ as $n\to\infty$.
\end{assumption}

%% ===================================================================
\section{Preliminary Lemmas}\label{sec:lemmas}
%% ===================================================================

%% ----- Oracle projection -----
\begin{lemma}[Oracle regression coefficient]\label{lem:oracle}
Under Assumption~\ref{ass:iso}, the best linear predictor of
$y^{(t)}$ given $x_{\cR(t)}$ (conditional on~$\beta^{(t)}$) is
$x_{\cR(t)}^\top\ostar$ with
\[
  \ostar = \Phi\,\beta^{(t)} = \beta^{(t)}_{\cR(t)}.
\]
The residual is
$y^{(t)}-x_{\cR(t)}^{\top}\ostar
 =x_{\cRc}^{\top}\beta^{(t)}_{\cRc}+\varepsilon$,
and is independent of~$x_{\cR(t)}$.
\end{lemma}

\begin{proof}
When $\Sigma=I_p$, the blocks satisfy
$\Sigma_{\cR(t)}=I_{p_t}$,
$\Sigma_{\cR(t),\cRc}=0$, and the Schur complement is
$\SigcR=I_{p-p_t}$.
Hence $\Phi=[I_{p_t}\mid 0]$, and
\[
  \ostar=\Phi\beta^{(t)}=\beta^{(t)}_{\cR(t)}.
\]
The conditional distribution
$x_{\cRc}\mid x_{\cR(t)}$ has mean $0$ and covariance $I_{p-p_t}$
(by the block--diagonal structure of $\Sigma=I_p$), so
$x_{\cRc}$ is independent of $x_{\cR(t)}$, and the residual
$x_{\cRc}^{\top}\beta^{(t)}_{\cRc}+\varepsilon$ is independent
of $x_{\cR(t)}$.
\end{proof}

%% ----- Closed form -----
\begin{lemma}[Closed form of the estimator]\label{lem:closedform}
For any $\lambda>0$,
\[
  \ohat
  =(\hat\Sigma+\lambda I)^{-1}
  \Bigl(\tfrac{1}{n}X^{\top}Y+\lambda\,\hat\omega_s^{(t)}\Bigr).
\]
\end{lemma}

\begin{proof}
The objective is strictly convex in~$\omega$.  Setting the gradient
to zero gives
\[
  -\tfrac{2}{n}X^{\top}(Y-X\omega)+2\lambda(\omega-\hat\omega_s^{(t)})=0
  \;\implies\;
  (\hat\Sigma+\lambda I)\omega=\tfrac{1}{n}X^{\top}Y+\lambda\,\hat\omega_s^{(t)}.
\]
Since $\hat\Sigma+\lambda I\succ 0$, the result follows.
\end{proof}

%% ----- Cross-term vanishing -----
\begin{lemma}[Vanishing of the cross term]\label{lem:cross}
Under Assumptions~\ref{ass:iso}--\ref{ass:indep}, write
$Y=X\ostar+\xi$ where $\xi_i=x_{i,\cRc}^{\top}\beta^{(t)}_{\cRc}+\varepsilon_i$.
Then for every deterministic matrix~$M$ that is a function of~$X$,
\[
  \E\!\bigl[\xi^{\top}\tfrac{1}{n}X\,M\,\Phi(\bhat-\beta^{(t)})
  \;\big|\;X\bigr]=0.
\]
\end{lemma}

\begin{proof}
Since $\Sigma=I_p$, $x_{i,\cRc}$ is independent of $x_{i,\cR(t)}$,
so $\xi$ is conditionally independent of $X$ given
$(\beta^{(t)}_{\cRc},\varepsilon)$.

Write $\xi=U\beta^{(t)}_{\cRc}+\varepsilon$ where $U\in\R^{n\times(p-p_t)}$
has rows $x_{i,\cRc}^{\top}$.  Since $U$ is independent of $X$, of
$\bhat$, and of $\varepsilon$:
\begin{align*}
  &\E[\xi^{\top}\tfrac{1}{n}X\,M\,\Phi(\bhat-\beta^{(t)})\mid X] \\
  &\quad
  =\E\bigl[
    \beta^{(t)\top}_{\cRc}
    \underbrace{\E[U^{\top}\mid X,\beta^{(t)},\bhat]}_{=\;0}
    \cdot\tfrac{1}{n}XM\Phi(\bhat-\beta^{(t)})
    \;\big|\;X\bigr] \\
  &\qquad
  +\E[\varepsilon^{\top}\tfrac{1}{n}XM\Phi(\bhat-\beta^{(t)})\mid X].
\end{align*}
The first term vanishes because $\E[U\mid X,\beta^{(t)},\bhat]=0$
(the rows of $U$ are mean-zero and independent of everything
conditioning on).
The second term vanishes because $\varepsilon$ is independent of
$(X,\bhat,\beta^{(t)})$.
\end{proof}

%% ----- Transfer error bound -----
\begin{lemma}[Bound on the transfer error covariance]\label{lem:delta}
Define $\delta=\Phi(\bhat-\beta^{(t)})\in\R^{p_t}$.  Under
Assumptions~\ref{ass:coef}--\ref{ass:iso},
\[
  \E[\delta\,\delta^{\top}]
  \;\preceq\;
  \frac{C_\delta}{p}\,I_{p_t},
  \qquad
  C_\delta\;:=\;\bigl(\sqrt{C^{U}}
  +\sqrt{\alpha_s^{2}+\alpha_t^{2}-2\rho\,\alpha_s\alpha_t}\;\bigr)^{2}.
\]
\end{lemma}

\begin{proof}
Write $\delta=\Phi(\bhat-\beta^{(s)})+\Phi(\beta^{(s)}-\beta^{(t)})
=:e_{\cR}+d_{\cR}$.

\medskip
\noindent\emph{Step 1 (Source estimation error).}\quad
Since $\Phi=[I_{p_t}\mid 0]$,
$e_{\cR}=(\bhat-\beta^{(s)})_{\cR(t)}$.
By the interlacing inequality for principal submatrices,
\[
  \lambda_{\max}\!\bigl(\E[e_{\cR}\,e_{\cR}^{\top}]\bigr)
  \;\le\;
  \lambda_{\max}\!\bigl(\E[(\bhat-\beta^{(s)})
  (\bhat-\beta^{(s)})^{\top}]\bigr)
  \;\le\;\frac{C^{U}}{p}.
\]

\noindent\emph{Step 2 (Coefficient shift).}\quad
$\beta^{(s)}_{\cR(t)}-\beta^{(t)}_{\cR(t)}$ has mean zero and
covariance
\[
  \E[d_{\cR}\,d_{\cR}^{\top}]
  =\frac{\alpha_s^{2}+\alpha_t^{2}-2\rho\,\alpha_s\alpha_t}{p}\,I_{p_t}
  =:\frac{\theta^{2}}{p}\,I_{p_t},
\]
using $\Cov(\beta^{(s)}-\beta^{(t)})
=(\alpha_s^{2}/p+\alpha_t^{2}/p-2\rho\alpha_s\alpha_t/p)\,I_p$.

\medskip
\noindent\emph{Step 3 (Cross term bound).}\quad
By the matrix Cauchy--Schwarz inequality (see, e.g., Bhatia~\cite{bhatia}),
\[
  \lambda_{\max}\!\bigl(\E[e_{\cR}\,d_{\cR}^{\top}]
  +\E[d_{\cR}\,e_{\cR}^{\top}]\bigr)
  \;\le\;
  2\,\lambda_{\max}(\E[e_{\cR}\,e_{\cR}^{\top}])^{1/2}\,
    \lambda_{\max}(\E[d_{\cR}\,d_{\cR}^{\top}])^{1/2}
  \;\le\;
  \frac{2\sqrt{C^{U}}\,\theta}{p}.
\]

\noindent\emph{Step 4 (Combining).}\quad
$\E[\delta\delta^{\top}]
=\E[e_{\cR}e_{\cR}^{\top}]+\E[d_{\cR}d_{\cR}^{\top}]
+\E[e_{\cR}d_{\cR}^{\top}]+\E[d_{\cR}e_{\cR}^{\top}]$, so
\[
  \lambda_{\max}\!\bigl(\E[\delta\delta^{\top}]\bigr)
  \;\le\;
  \frac{C^{U}+\theta^{2}+2\sqrt{C^{U}}\,\theta}{p}
  =\frac{(\sqrt{C^{U}}+\theta)^{2}}{p}
  =\frac{C_\delta}{p},
\]
whence $\E[\delta\delta^{\top}]\preceq(C_\delta/p)\,I_{p_t}$.
\end{proof}

%% ----- RMT -----
\begin{lemma}[Marchenko--Pastur trace functionals]%
\label{lem:MP}
Suppose $X\in\R^{n\times p_t}$ has i.i.d.\ rows drawn from
$\N(0,I_{p_t})$ and $p_t/n\to\eta\in(0,\infty)$.  For fixed
$\lambda>0$, define
\[
  \phi_1(\lambda,\eta)
  =\lim_{n\to\infty}\frac{1}{p_t}
   \tr\!\bigl(\hat\Sigma(\hat\Sigma+\lambda I)^{-2}\bigr),
  \qquad
  \phi_2(\lambda,\eta)
  =\lim_{n\to\infty}\frac{1}{p_t}
   \tr\!\bigl((\hat\Sigma+\lambda I)^{-2}\bigr).
\]
Both limits exist almost surely.  They are determined by the
Marchenko--Pastur law $F_\eta$ with density
\[
  dF_\eta(t)=\frac{1}{2\pi\eta t}
  \sqrt{(b-t)(t-a)}\;\mathbf{1}_{[a,b]}(t)\,dt
  +\Bigl(1-\frac{1}{\eta}\Bigr)^{+}\!\delta_0(t),
  \qquad a=(1{-}\sqrt\eta)^{2},\;b=(1{+}\sqrt\eta)^{2},
\]
via
\begin{equation}\label{eq:phi12}
  \phi_1(\lambda,\eta)
  =\int_{0}^{\infty}\frac{t}{(t+\lambda)^{2}}\,dF_\eta(t),
  \qquad
  \phi_2(\lambda,\eta)
  =\int_{0}^{\infty}\frac{1}{(t+\lambda)^{2}}\,dF_\eta(t).
\end{equation}
Moreover, $\phi_1+\lambda\,\phi_2 = m_\eta(-\lambda)$, the Stieltjes
transform of~$F_\eta$ evaluated at~$-\lambda$, which satisfies
the self-consistent equation
\begin{equation}\label{eq:stieltjes}
  m_\eta(-\lambda)
  =\int\frac{dF_\eta(t)}{t+\lambda}
  =\frac{1-\eta+\lambda
  -\sqrt{(1-\eta+\lambda)^{2}+4\eta\lambda}}{2\eta\lambda},
  \quad\lambda>0.
\end{equation}
\end{lemma}

\begin{proof}
By the classical Marchenko--Pastur theorem~\cite{MP1967}, the
empirical spectral distribution of $\hat\Sigma$ converges to
$F_\eta$.  The almost sure convergence of the trace functionals
follows from the dominated convergence theorem applied to the a.s.\
converging spectral measure, together with the uniform integrability
provided by the fact that
$t/(t+\lambda)^{2}\le 1/(4\lambda)$ and $1/(t+\lambda)^{2}\le
1/\lambda^{2}$ are bounded continuous functions for $\lambda>0$.
The closed form \eqref{eq:stieltjes} is classical; see
\cite{silverstein1995strong}.
\end{proof}

%% ===================================================================
\section{Main Result}\label{sec:main}
%% ===================================================================

\begin{theorem}[Asymptotic predictive risk upper bound]%
\label{thm:main}
Under Assumptions~\ref{ass:dim}--\ref{ass:lambda}, for every fixed
$\lambda>0$ the predictive risk satisfies
\begin{equation}\label{eq:riskbound}
\boxed{
  \limsup_{n,p\to\infty}
  \E\!\left[r^{\mathrm{pred}}(\ohat)\right]
  \;\le\;
  \mathcal R_{\eta}\!\left(\lambda;C_\delta,\sigma_\xi^2\right),
}
\end{equation}
where $\mathcal R_{\eta}(\lambda;a^2,\tau^2)$ denotes the
Marchenko--Pastur deterministic-equivalent predictive-risk functional
for isotropic ridge regression with aspect ratio $\eta$, signal
strength $a^2$, and noise variance $\tau^2$.  Here
\begin{gather*}
  \sigma_\xi^{2}=\alpha_t^{2}\Bigl(1-\frac{\eta}{\gamma}\Bigr)
  +\sigma^{2},
  \qquad
  \theta^{2}=\alpha_s^{2}+\alpha_t^{2}-2\rho\,\alpha_s\alpha_t,
  \qquad
  C_\delta=(\sqrt{C^{U}}+\theta)^{2},
\end{gather*}

Consequently, the optimally tuned upper bound satisfies
\begin{equation}\label{eq:lambdastar}
  \lambda^{\star}=
  \frac{\sigma_\xi^{2}\,\gamma}{C_\delta}.
\end{equation}
Equivalently,
\begin{equation}\label{eq:main_unexpanded_opt}
\boxed{
  \inf_{\lambda>0}\limsup_{n,p\to\infty}
  \E\!\left[r^{\mathrm{pred}}(\ohat)\right]
  \;\le\;
  \mathcal R_{\eta}\!\left(\lambda^\star;C_\delta,\sigma_\xi^2\right).
}
\end{equation}
\end{theorem}

%% ===================================================================
\section{Proof of Theorem~\ref{thm:main}}\label{sec:proof}
%% ===================================================================

The proof proceeds in five steps.

%% ----- Step 1 -----
\subsection*{Step 1: Bias--variance--missing feature decomposition}

For a new test point $(\tilde x,\tilde y^{(t)})$ from the target
site, write
\[
  \tilde y^{(t)}-\tilde x_{\cR(t)}^{\top}\ohat
  =\underbrace{\tilde x_{\cR(t)}^{\top}(\ostar-\ohat)}_{\text{(I)}}
  +\underbrace{\tilde x_{\cRc}^{\top}\beta^{(t)}_{\cRc}}_{\text{(II)}}
  +\underbrace{\tilde\varepsilon}_{\text{(III)}}.
\]

By Lemma~\ref{lem:oracle}, (II) is independent of
$\tilde x_{\cR(t)}$, and (III) is independent of everything
else.  Squaring and taking $\E[\,\cdot\mid X]$:
\begin{itemize}[nosep]
\item The cross term between (I) and (III) vanishes because
$\E[\tilde\varepsilon]=0$ and $\tilde\varepsilon$ is independent of
everything else.
\item The cross term between (I) and (II): (II) depends on
$\tilde x_{\cRc}$ and $\beta^{(t)}_{\cRc}$; (I) depends on
$\tilde x_{\cR(t)}$ and the training data.  Conditioning first on all
randomness except $\tilde x_{\cRc}$, we get
$\E[\tilde x_{\cRc}\mid\tilde x_{\cR(t)}]=0$ (by $\Sigma=I_p$), so
this cross term vanishes.
\item The cross term between (II) and (III) vanishes since
$\tilde\varepsilon\perp \tilde x$.
\end{itemize}

Therefore
\begin{equation}\label{eq:decomp}
  r^{\mathrm{pred}}(\ohat)
  =\underbrace{\E\bigl[
  (\ostar-\ohat)^{\top}\SigR(\ostar-\ohat)
  \;\big|\;X\bigr]}_{R_{\mathrm{est}}}
  +\underbrace{\E\bigl[
  \beta^{(t)\top}_{\cRc}\SigcR\,\beta^{(t)}_{\cRc}
  \bigr]}_{R_{\mathrm{miss}}}
  +\;\sigma^{2},
\end{equation}
where we used
$\E[\tilde x_{\cR(t)}\tilde x_{\cR(t)}^{\top}]=\SigR=I_{p_t}$
and $\SigcR=I_{p-p_t}$.

The missing-feature risk is
\[
  R_{\mathrm{miss}}
  =\tr\!\bigl(I_{p-p_t}\cdot\tfrac{\alpha_t^{2}}{p}\,I_{p-p_t}\bigr)
  =\frac{\alpha_t^{2}(p-p_t)}{p}
  =\alpha_t^{2}\Bigl(1-\frac{p_t}{p}\Bigr)
  \;\xrightarrow{\;n\to\infty\;}
  \;\alpha_t^{2}\Bigl(1-\frac{\eta}{\gamma}\Bigr).
\]

%% ----- Step 2 -----
\subsection*{Step 2: Estimation error decomposition}

By Lemma~\ref{lem:closedform} and the identity $Y=X\ostar+\xi$
(Lemma~\ref{lem:oracle}), set $M=(\hat\Sigma+\lambda I)^{-1}$ and
compute
\begin{align}
  \ohat-\ostar
  &=M\bigl(\hat\Sigma\,\ostar+\tfrac{1}{n}X^{\top}\xi
  +\lambda\,\hat\omega_s^{(t)}\bigr)-\ostar
  \notag\\
  &=M\bigl(\tfrac{1}{n}X^{\top}\xi
  +\lambda(\hat\omega_s^{(t)}-\ostar)\bigr)
  \notag\\
  &=\underbrace{M\,\tfrac{1}{n}X^{\top}\xi}_{A}
  \;+\;\underbrace{\lambda\,M\,\delta}_{B},
  \label{eq:error_decomp}
\end{align}
where $\delta=\Phi(\bhat-\beta^{(t)})
=\hat\beta^{(s)}_{\cR(t)}-\beta^{(t)}_{\cR(t)}$.

Since $\SigR=I_{p_t}$,
\[
  R_{\mathrm{est}}
  =\E[\|A+B\|^{2}\mid X]
  =\E[\|A\|^{2}\mid X]
  +\E[\|B\|^{2}\mid X]
  +2\,\E[A^{\top}B\mid X].
\]

By Lemma~\ref{lem:cross}, the cross term
$\E[A^{\top}B\mid X]=0$.  Hence
\begin{equation}\label{eq:Rest}
  R_{\mathrm{est}}
  =\underbrace{\E[\|A\|^{2}\mid X]}_{T_1(\text{variance})}
  +\underbrace{\E[\|B\|^{2}\mid X]}_{T_2(\text{bias})}.
\end{equation}

%% ----- Step 3 -----
\subsection*{Step 3: Evaluation of the variance term $T_1$}

We compute $\E[\xi\xi^{\top}\mid X]$.  Since $\Sigma=I_p$, the rows
of $U$ (the matrix with rows $x_{i,\cRc}^{\top}$) are i.i.d.\
$\N(0,I_{p-p_t})$ and \emph{independent} of~$X$.  With $\varepsilon$
also independent:
\[
  \E[\xi_i\xi_j\mid X]
  =\begin{cases}
  \E\bigl[(x_{i,\cRc}^{\top}\beta^{(t)}_{\cRc})^{2}\bigr]+\sigma^{2}
  &\text{if }i=j,\\
  0 &\text{if }i\neq j,
  \end{cases}
\]
since $\E[x_{i,\cRc}]=0$ independently across $i$, and
$\varepsilon_i\perp\varepsilon_j$ for $i\neq j$.

For the diagonal entry:
\[
  \E\bigl[(x_{i,\cRc}^{\top}\beta^{(t)}_{\cRc})^{2}\bigr]
  =\E\bigl[\beta^{(t)\top}_{\cRc}
  \underbrace{\E[x_{i,\cRc}x_{i,\cRc}^{\top}]}_{I_{p-p_t}}
  \beta^{(t)}_{\cRc}\bigr]
  =\E[\|\beta^{(t)}_{\cRc}\|^{2}]
  =\frac{\alpha_t^{2}(p-p_t)}{p}.
\]

Define
$\sigma_\xi^{2}:=\alpha_t^{2}(p-p_t)/p+\sigma^{2}$.  Then
$\E[\xi\xi^{\top}\mid X]=\sigma_\xi^{2}\,I_n$, and
\begin{align}
  T_1
  &=\E\bigl[\tfrac{1}{n^{2}}\xi^{\top}XM^{2}X^{\top}\xi\;\big|\;X\bigr]
  =\frac{\sigma_\xi^{2}}{n^{2}}\tr(XM^{2}X^{\top})
  =\frac{\sigma_\xi^{2}}{n}\tr(M^{2}\hat\Sigma)
  \notag\\
  &=\frac{\sigma_\xi^{2}}{n}\tr\!\bigl(
  \hat\Sigma(\hat\Sigma+\lambda I)^{-2}\bigr)
  =\sigma_\xi^{2}\,\frac{p_t}{n}\cdot
  \frac{1}{p_t}\tr\!\bigl(\hat\Sigma(\hat\Sigma+\lambda I)^{-2}\bigr).
  \label{eq:T1}
\end{align}

By Lemma~\ref{lem:MP},
$\frac{1}{p_t}\tr(\hat\Sigma(\hat\Sigma+\lambda I)^{-2})
\xrightarrow{a.s.}\phi_1(\lambda,\eta)$, so
\begin{equation}\label{eq:T1asymp}
  T_1\;\xrightarrow{a.s.}\;
  \sigma_\xi^{2}\,\eta\,\phi_1(\lambda,\eta).
\end{equation}

%% ----- Step 4 -----
\subsection*{Step 4: Upper bound on the bias term $T_2$}

\begin{align}
  T_2
  &=\lambda^{2}\,\E[\delta^{\top}M^{2}\delta\mid X]
  =\lambda^{2}\,\tr\!\bigl(M^{2}\,\E[\delta\delta^{\top}]\bigr).
  \label{eq:T2trace}
\end{align}
(Here $M$ is a deterministic function of $X$, and $\delta$ is
independent of $X$ only through the residual randomness of $\beta$
and $\bhat$.  The expectation is over all remaining randomness.)

By Lemma~\ref{lem:delta},
$\E[\delta\delta^{\top}]\preceq(C_\delta/p)\,I_{p_t}$, hence
\begin{align}
  T_2
  &\le\frac{\lambda^{2}\,C_\delta}{p}\,
  \tr\!\bigl((\hat\Sigma+\lambda I)^{-2}\bigr)
  =\frac{\lambda^{2}\,C_\delta\,p_t}{p}\cdot
  \frac{1}{p_t}\tr\!\bigl((\hat\Sigma+\lambda I)^{-2}\bigr).
  \label{eq:T2bound}
\end{align}

By Lemma~\ref{lem:MP},
$\frac{1}{p_t}\tr((\hat\Sigma+\lambda I)^{-2})
\xrightarrow{a.s.}\phi_2(\lambda,\eta)$, and
$p_t/p\to\eta/\gamma$, so
\begin{equation}\label{eq:T2asymp}
  T_2\;\le\;
  \frac{\lambda^{2}\,C_\delta\,\eta}{\gamma}\,\phi_2(\lambda,\eta)
  +o(1).
\end{equation}

%% ----- Step 5 -----
\subsection*{Step 5: Combining and optimising}

Substituting \eqref{eq:T1asymp} and \eqref{eq:T2asymp} into
\eqref{eq:decomp}--\eqref{eq:Rest}, and noting that
$\sigma_\xi^{2}\to\alpha_t^{2}(1-\eta/\gamma)+\sigma^{2}$, we obtain
the claimed bound~\eqref{eq:riskbound}.

\medskip
\noindent\textbf{Optimisation over $\lambda$.}\quad
Define the estimation-error upper bound
\[
  \overline{R}_{\mathrm{est}}(\lambda)
  =\sigma_\xi^{2}\,\eta\,\phi_1(\lambda,\eta)
  +\frac{\lambda^{2}\,C_\delta\,\eta}{\gamma}\,\phi_2(\lambda,\eta).
\]
Setting $\partial\overline{R}_{\mathrm{est}}/\partial\lambda=0$:
\begin{equation}\label{eq:FOC}
  \sigma_\xi^{2}\,\eta\,\phi_1'(\lambda,\eta)
  +\frac{C_\delta\,\eta}{\gamma}\bigl(
  2\lambda\,\phi_2(\lambda,\eta)+\lambda^{2}\phi_2'(\lambda,\eta)
  \bigr)=0.
\end{equation}
The identity
$\phi_1(\lambda,\eta)=\int t(t+\lambda)^{-2}\,dF_\eta$ yields
$\phi_1'(\lambda,\eta)=-2\int t(t+\lambda)^{-3}\,dF_\eta$, and similarly
$\phi_2'(\lambda,\eta)=-2\int(t+\lambda)^{-3}\,dF_\eta$.
Hence $\phi_1'=-2\int t(t+\lambda)^{-3}$ and
$2\phi_2+\lambda\phi_2'
=2\int(t+\lambda)^{-2}-2\lambda\int(t+\lambda)^{-3}
=2\int t(t+\lambda)^{-3}=-\phi_1'$.

Substituting into \eqref{eq:FOC}:
\[
  \sigma_\xi^{2}\,\eta\,\phi_1'
  +\frac{C_\delta\,\eta\,\lambda}{\gamma}(-\phi_1')=0
  \;\implies\;
  \phi_1'\Bigl(\sigma_\xi^{2}-\frac{C_\delta\,\lambda}{\gamma}\Bigr)=0.
\]
Since $\phi_1'(\lambda,\eta)<0$ for all $\lambda>0$ (the integrand is
strictly positive on the support of~$F_\eta$), we conclude
\begin{equation}\label{eq:lambdastar_simple}
  \boxed{\lambda^{\star}=\frac{\sigma_\xi^{2}\,\gamma}{C_\delta}.}
\end{equation}

\medskip
\noindent\textbf{Optimal risk.}\quad
Substituting $\lambda^{\star}$ into
$\overline{R}_{\mathrm{est}}$:
\begin{align}
  \overline{R}_{\mathrm{est}}(\lambda^{\star})
  &=\sigma_\xi^{2}\,\eta\,\phi_1(\lambda^{\star},\eta)
  +\frac{(\sigma_\xi^{2}\gamma/C_\delta)^{2}\,C_\delta\,\eta}
  {\gamma}\,\phi_2(\lambda^{\star},\eta)
  \notag\\
  &=\sigma_\xi^{2}\,\eta\bigl[
  \phi_1(\lambda^{\star},\eta)
  +\lambda^{\star}\,\phi_2(\lambda^{\star},\eta)\bigr]
  \notag\\
  &=\sigma_\xi^{2}\,\eta\,m_\eta(-\lambda^{\star}),
  \label{eq:optRest}
\end{align}
where the last equality uses
$\phi_1+\lambda\phi_2=\int\frac{t+\lambda}{(t+\lambda)^{2}}\,dF_\eta
=\int\frac{1}{t+\lambda}\,dF_\eta=m_\eta(-\lambda)$.

Thus the optimally tuned predictive-risk upper bound can be reported in
the unexpanded deterministic-equivalent form
\begin{equation}\label{eq:optbound}
  \boxed{
  \inf_{\lambda>0}\limsup_{n,p\to\infty}
  \E\!\left[r^{\mathrm{pred}}(\ohat)\right]
  \;\le\;
  \mathcal R_{\eta}\!\left(\lambda^\star;C_\delta,\sigma_\xi^2\right),
  \qquad
  \lambda^{\star}=\frac{\sigma_\xi^{2}\,\gamma}{C_\delta}.
  }
\end{equation}
This completes the proof.

%% ===================================================================
\section{Remarks}\label{sec:remarks}
%% ===================================================================

\begin{remark}[Interpretation of $\lambda^{\star}$]\label{rem:lambda}
The optimal regularisation
$\lambda^{\star}=\sigma_\xi^{2}\gamma/C_\delta$ balances the
effective noise level $\sigma_\xi^{2}$ (which includes
the irreducible error from missing features) against the transfer
error~$C_\delta$.  When the source estimator is very
accurate ($C^{U}\to 0$) and the sites are nearly identical
($\rho\to 1$, $\alpha_s=\alpha_t$), $C_\delta\to 0$, so
$\lambda^{\star}\to\infty$: we rely almost entirely on the
transferred estimate.  Conversely, when the source information is
poor ($C_\delta$ large), $\lambda^{\star}\to 0$ and the estimator
reduces to the ordinary least squares (or Moore--Penrose) solution
on the target data alone.
\end{remark}

\begin{remark}[Comparison with no transfer]\label{rem:notransfer}
Without transfer ($\lambda=0$, provided $\eta<1$ so that
$\hat\Sigma$ is invertible), the estimation error is
$\sigma_\xi^{2}\eta\phi_1(0,\eta)$.  Since
$\phi_1(0,\eta)=\int t^{-1}\,dF_\eta(t)$, which diverges as
$\eta\nearrow 1$, the transfer estimator with optimal
$\lambda^{\star}>0$ provides a strict improvement whenever
$\eta$ is not too small and $C_\delta$ is finite.  When $\eta\ge 1$,
the no-transfer estimator is not even defined (the OLS problem is
underdetermined), while the transfer estimator remains well-defined
for any $\lambda>0$.
\end{remark}

\begin{remark}[General covariance $\Sigma$]\label{rem:general_Sigma}
When $\Sigma\neq I_p$, the analysis extends with the following
modifications:
(i)~the oracle coefficient becomes
$\ostar=\beta^{(t)}_{\cR(t)}
+\SigR^{-1}\SigRRc\,\beta^{(t)}_{\cRc}=\Phi\beta^{(t)}$;
(ii)~the residual covariance becomes
$\sigma_\xi^{2}=(\alpha_t^{2}/p)\,\tr(\SigcR)+\sigma^{2}$;
(iii)~the trace functionals $\phi_1,\phi_2$ are replaced by
their generalisations for the generalised Marchenko--Pastur law
associated with $\SigR$; and
(iv)~the transfer error bound requires
$\E[\delta\delta^{\top}]\preceq(C_\delta/p)\,\Phi\Phi^{\top}$,
which involves $\|\Phi\|_{\mathrm{op}}$ and the condition number
of $\SigR$.  The structure of the proof is identical; the additional
spectral assumptions on $\Sigma$ (bounded eigenvalues, bounded away
from~$0$) ensure the random-matrix-theory limits exist.
\end{remark}

\begin{remark}[Sharpness of the bound]\label{rem:sharp}
The bound is an \emph{equality} (up to the $o(1)$ term) under two
conditions:
(a)~the cross terms in Lemma~\ref{lem:delta} vanish, i.e.\
$\E[e_{\cR}\,d_{\cR}^{\top}]=0$ (which holds, for example, when
$\bhat$ is an unbiased estimator of $\beta^{(s)}$ constructed
independently of the random coefficients), and
(b)~$\E[\delta\delta^{\top}]=(C_\delta/p)\,I_{p_t}$, which holds
when $\E[e_{\cR}e_{\cR}^{\top}]=(C^{U}/p)\,I_{p_t}$ exactly
(e.g.\ when the source estimator is a ridge estimator with isotropic
error covariance).
In that case the bound~\eqref{eq:optbound} is the exact asymptotic
risk.
\end{remark}

%% ===================================================================
\begin{thebibliography}{9}
\bibitem{MP1967}
V.~A.~Mar\v{c}enko and L.~A.~Pastur,
``Distribution of eigenvalues for some sets of random matrices,''
\emph{Math.\ USSR-Sb.}, vol.~1, no.~4, pp.~457--483, 1967.

\bibitem{silverstein1995strong}
J.~W.~Silverstein and Z.~D.~Bai,
``On the empirical distribution of eigenvalues of a class of large
dimensional random matrices,''
\emph{J.\ Multivariate Anal.}, vol.~54, pp.~175--192, 1995.

\bibitem{bhatia}
R.~Bhatia,
\emph{Matrix Analysis},
Springer, 1997.

\bibitem{dobriban2018}
E.~Dobriban and S.~Wager,
``High-dimensional asymptotics of prediction: Ridge regression and
classification,''
\emph{Ann.\ Statist.}, vol.~46, no.~1, pp.~247--279, 2018.

\bibitem{hastie2022surprises}
T.~Hastie, A.~Montanari, S.~Rosset, and R.~J.~Tibshirani,
``Surprises in high-dimensional ridgeless least squares interpolation,''
\emph{Ann.\ Statist.}, vol.~50, no.~2, pp.~949--986, 2022.
\end{thebibliography}

\end{document}
